\begin{table}[t]
\centering
\caption{Standardized Effect Sizes}
\label{tab:sde}
\footnotesize
\begin{tabular}{lcccccc}
\hline\hline
Outcome & $\hat{\beta}$ & SE & SD($Y$) & SDE & SE(SDE) & Classification \\
\hline
\multicolumn{7}{l}{\textit{Panel A: Pooled (2019--2023)}} \\[3pt]
Property crimes (DSP) & 0.198 & 0.327 & 45.4 & 0.0220 & 0.0363 & Small positive \\
Property crimes (payday) & 3.980 & 2.239 & 45.4 & 0.0876 & 0.0493 & Moderate positive \\
Robbery (DSP) & 0.013 & 0.175 & 26.4 & 0.0025 & 0.0335 & Null \\
Theft (DSP) & 0.185 & 0.165 & 25.4 & 0.0367 & 0.0328 & Small positive \\
{}[6pt]
\multicolumn{7}{l}{\textit{Panel B: Heterogeneous (period splits)}} \\[3pt]
Property crimes, 2019 & 0.625 & 0.473 & 45.4 & 0.0694 & 0.0526 & Moderate positive \\
Property crimes, 2022--23 & -0.424 & 0.340 & 54.7 & -0.0391 & 0.0314 & Small negative \\
\hline
\end{tabular}
\par\vspace{0.5em}
{\scriptsize \textit{Notes:} \textbf{Country:} Argentina. \textbf{Research question:} Does welfare payment timing affect property crime through a liquidity depletion cycle in Buenos Aires? \textbf{Policy mechanism:} ANSES distributes monthly pension and social transfer payments to approximately 18 million beneficiaries on staggered calendar days determined by the last digit of each person's national identity document (DNI), creating quasi-random within-month variation in neighborhood-level cash availability. \textbf{Outcome definition:} Daily count of property crimes (robbery plus theft) recorded by Buenos Aires City's Ministry of Justice and Security. \textbf{Treatment:} Continuous---average number of days since most recent ANSES payment across all 10 DNI digit groups; and binary---indicator for whether at least one digit group received payment on a given day. \textbf{Data:} Buenos Aires City Open Data crime records and ANSES payment calendar, 2019--2023, city-day observations, $N$ = 1,819 days. \textbf{Method:} OLS with day-of-week and year$\times$month fixed effects, standard errors clustered by year$\times$month. \textbf{Sample:} All days 2019--2023 with complete payment calendar data; Panel B splits by pre-COVID (2019) and post-COVID (2022--2023) periods. SDE $= \hat{\beta} \times \text{SD}(X) / \text{SD}(Y)$ for continuous treatment; SDE $= \hat{\beta} / \text{SD}(Y)$ for binary treatment, where SD($Y$) is the pre-treatment (2019) standard deviation. Classification refers to magnitude, not statistical significance: Large ($|$SDE$| > 0.15$), Moderate ($0.05$--$0.15$), Small ($0.005$--$0.05$), Null ($< 0.005$).}
\end{table}
